\chapter*{Abstract}\label{cha:abstract}

Top-fermented beers owe much of their fruity aroma to esters and higher alcohols formed by the
yeast during fermentation. This (fictional) thesis investigates how strongly the fermentation
temperature shapes this aroma profile. In twelve brews of a pale ale the temperature was varied
between 16\,\textdegree{}C and 24\,\textdegree{}C. Isoamyl acetate -- the typical banana note --
increased almost linearly by about 38\,\% per 4\,K. A blind tasting with
24 participants confirms: warmer fermented beers are perceived as fruitier, but less ``clean''.

\chapter*{Kurzfassung}\label{cha:abstract_german}

Obergärige Biere verdanken ihr fruchtiges Aroma vor allem den Estern und
höheren Alkoholen, die Hefen während der Gärung bilden. Diese (erfundene) Arbeit untersucht, wie
stark die Gärtemperatur dieses Aromaprofil prägt. Isoamylacetat stieg um rund 38\,\% je
4\,K; eine Blindverkostung bestätigt den Effekt.
